% Computer Science A --- Data structures, collections, and algorithm basics (\input from \texttt{main.tex}).
% Free-response bank at the front is migrated from \texttt{cs-a/05-practice.tex} (Collections Framework, DSA basics). Design-pattern FRQs live in \texttt{object\_oriented.tex}.

\runningheader{Computer Science A}{Data structures \& algorithms}{Page \thepage\ of \numpages}

\begingradingrange{csau5}

\uplevel{\par\textbf{The Java Collections Framework}\par}

\question[4]
What is a \texttt{Collection}?
\begin{solutionorbox}[1.15in]
  Root interface in the Java Collections Framework for groups of objects (\texttt{add}, \texttt{remove}, \texttt{size}, iteration); \texttt{Map} is not a \texttt{Collection} but is part of the same library.
\end{solutionorbox}

\question[4]
What is the Collection API?
\begin{solutionorbox}[1.2in]
  The set of interfaces and implementations in \texttt{java.util} (and related) for lists, sets, queues, maps, sorting, and utilities such as \texttt{Collections} and \texttt{Arrays}.
\end{solutionorbox}

\question[4]
What are the major differences between a \texttt{List}, \texttt{Queue}, \texttt{Set}, and \texttt{Map}?
\begin{solutionorbox}[1.65in]
  \texttt{List}: ordered, indexed, duplicates allowed. \texttt{Queue}: FIFO (or priority) processing order. \texttt{Set}: no duplicates, often unordered (\texttt{TreeSet} sorted). \texttt{Map}: key--value pairs, keys unique; not a \texttt{Collection} of entries in the same sense.
\end{solutionorbox}

\question[4]
What is the difference between a \texttt{Stack} and a \texttt{Queue}?
\begin{solutionorbox}[1.2in]
  Stack is LIFO (last in, first out); queue is FIFO (first in, first out). Prefer \texttt{Deque} over legacy \texttt{Stack} in modern Java.
\end{solutionorbox}

\newpage

\uplevel{\par\textbf{Data structures and algorithms}\par}

\question[4]
Describe a linked list.
\begin{solutionorbox}[1.35in]
  Nodes where each holds data and a reference to the next (and previous for doubly linked); supports \(O(1)\) insert/delete at a known node but \(O(n)\) indexed access unless optimized.
\end{solutionorbox}

\question[4]
What is an algorithm?
\begin{solutionorbox}[1.05in]
  A finite, well-defined procedure that transforms input into output; analyzable for correctness and resource use.
\end{solutionorbox}

\question[4]
How would you describe time complexity?
\begin{solutionorbox}[1.35in]
  Big-\(O\) (and related) notation for how runtime grows with input size \(n\), ignoring constants and lower-order terms; describes worst/average cases as appropriate.
\end{solutionorbox}

\question[4]
What is the difference between linear search and binary search?
\begin{solutionorbox}[1.35in]
  Linear scans each element in order---\(O(n)\). Binary search halves a sorted range each step---\(O(\log n)\) but requires sorted (or indexed) data.
\end{solutionorbox}

\question[4]
What is the difference between Bubble Sort and Merge Sort?
\begin{solutionorbox}[1.45in]
  Bubble sort repeatedly swaps adjacent out-of-order elements; simple but typically \(O(n^2)\). Merge sort divides and merges sorted halves; \(O(n\log n)\) time, uses extra space for merging.
\end{solutionorbox}

\endgradingrange{csau5}

\uplevel{\par\textbf{Collections, casting, and wrapper types}\par}

\question[4]
What is the difference between an array and an \texttt{ArrayList}?
\begin{solutionorbox}[1.45in]
  Arrays have fixed length and can hold primitives or references; \texttt{ArrayList} is a resizable list implementation of \texttt{List} holding references (wrappers for primitives), with convenient add/remove and growth.
\end{solutionorbox}

\question[4]
What is explicit and implicit casting?
\begin{solutionorbox}[1.25in]
  Implicit casting is an automatic widening conversion (e.g.\ \texttt{int} to \texttt{long}) when safe; explicit casting uses parentheses and may narrow or cross types (e.g.\ \texttt{(int) d}), possibly losing precision or requiring runtime checks for references.
\end{solutionorbox}

\question[4]
What is autoboxing and unboxing?
\begin{solutionorbox}[1.2in]
  Compiler-inserted conversion between primitives and their wrapper types (e.g.\ \texttt{int} \(\leftrightarrow\) \texttt{Integer}) when assigning or passing to generic collections expecting objects.
\end{solutionorbox}

\question[4]
What are some operations you can perform on a \texttt{List}?
\begin{solutionorbox}[1.35in]
  Add/remove, get/set by index, \texttt{size}, iteration, \texttt{contains}/\texttt{indexOf}, sublists, sort, clear---details depend on the concrete \texttt{List} class.
\end{solutionorbox}

\newpage

\begingradingrange{csadsa}

\question[4]
Fill in the code from given choices.
\begin{lstlisting}[style=java]
import java.util.*;
public class Main {
    public static void main(String[] args) {
        ArrayList<Integer> al = new ArrayList<>();
        al.add(1);
        al.add(13);
        al.add(2);
        al.add(0);
        //fill here
        {
            System.out.println(i);
        }
    }
}
\end{lstlisting}
\begin{choices}
  \CorrectChoice \texttt{for(Integer i:al)}
  \choice \texttt{for(String i: al)}
  \choice \texttt{foreach(int i:al)}
  \choice \texttt{foreach(Integer i:al)}
\end{choices}
\begin{solution}
  Java's enhanced \texttt{for} loop (``for-each'') iterates over \texttt{al} with element type \texttt{Integer}. The loop variable type must match the collection's element type; \texttt{String} is wrong. Java has no \texttt{foreach} keyword---use \texttt{for} with a type and identifier before the colon.
\end{solution}

\question[4]
The methods that can be used by an iterator when traversing a collection are\ldots?
\begin{choices}
  \choice \texttt{hasNext()}
  \choice \texttt{next()}
  \choice \texttt{remove()}
  \CorrectChoice All of the above
\end{choices}
\begin{solution}
  \texttt{Iterator} supports \texttt{hasNext()} and \texttt{next()} on all collections that return an iterator. \texttt{remove()} is optional in the sense that it removes the last element returned by \texttt{next()} and may throw \texttt{UnsupportedOperationException} for read-only views, but it is part of the \texttt{Iterator} API for mutable traversals.
\end{solution}

\question[4]
True or False: Because sets have no ordering, there is no way to retrieve a specific object from a Set immediately.
\begin{choices}
  \CorrectChoice TRUE
  \choice FALSE
\end{choices}
\begin{solution}
  \texttt{Set} has no index-based access like \texttt{list.get(i)}. You can test membership with \texttt{contains} or remove by value, but you cannot ``fetch the $k$-th element'' or random-access a chosen slot the way you can with a \texttt{List}. The statement is interpreted in that sense (no immediate positional retrieval).
\end{solution}


\newpage

\question[4]
``Collections'' is?
\begin{choices}
  \choice The interface that all java collections implement
  \CorrectChoice A class filled with static methods used to manipulate collections
\end{choices}
\begin{solution}
  \texttt{java.util.Collections} is a \textbf{utility class} of static helpers: sort, binary search, unmodifiable wrappers, synchronized wrappers, etc. The root interface for many collection types is \texttt{Collection}; \texttt{Collections} (with an \texttt{s}) is not that interface.
\end{solution}

\question[4]
True or False: An array is a type of \texttt{Collection} in Java.
\begin{choices}
  \choice TRUE
  \CorrectChoice FALSE
\end{choices}
\begin{solution}
  Arrays are objects with \texttt{length} and index syntax \texttt{[ ]}, but they do not implement \texttt{Collection}. Use \texttt{Arrays.asList(\ldots)} or copy into an \texttt{ArrayList} when you need the \texttt{Collection} API.
\end{solution}

\question[4]
What should be the output?
\begin{lstlisting}[style=java,numbers=left]
import java.util.*;
class SetTest
{
  public static void main(String args[])
  {
    Set<Object> s = new HashSet<Object>();

    s.add("java");
    s.add(new String("java"));
    System.out.println(s);
  }
}
\end{lstlisting}
\begin{choices}
  \choice Code will compile and print java twice
  \choice Code will not compile as duplicates in set are not allowed
  \CorrectChoice Code will compile and put java as output
  \choice Code will not compile because of line 8,9 as in set of type object, we cannot add string objects
\end{choices}
\begin{solution}
  \texttt{String} equality uses value equality: \texttt{"java"} and \texttt{new String("java")} are \texttt{equals}, so the set holds \textbf{one} element. \texttt{HashSet} allows \texttt{String} instances in an \texttt{Set<Object>}. Output is typically a single-line set representation such as \texttt{[java]} (order not guaranteed).
\end{solution}


\newpage

\question[4]
What should be the output of this code snippet?
\begin{lstlisting}[style=java]
import java.util.ArrayList;
import java.util.List;
public class MyClass {
    public static void main(String[] args) {
        List<Integer> list = new ArrayList<Integer>();
        list.add(2);
        list.add(3);
        m(list);
    }
    public static void m(List<Number> list) {
        System.out.println(list);
    }
}
\end{lstlisting}
\begin{choices}
  \choice Prints \texttt{[2,3]}
  \choice Runtime Exception
  \CorrectChoice Compile time error
  \choice None of the above
\end{choices}
\begin{solution}
  \texttt{List<Integer>} is \textbf{not} a subtype of \texttt{List<Number>} in Java (generics are invariant). Passing \texttt{list} into \texttt{m} is a compile-time error. Intuitively, if it compiled, \texttt{m} could add a \texttt{Double} into a \texttt{List<Integer>}.
\end{solution}

\question[4]
What are Java \textbf{generics}?
\begin{choices}
  \choice They reduce the amount of compile-time checking
  \choice They add capabilities when defining \emph{only} primitive variables
  \CorrectChoice They let types (classes and interfaces) be parameters when defining classes, interfaces, and methods
  \choice They give collections the ability to catch exceptions automatically
\end{choices}
\begin{solution}
  Generics parameterize types (\texttt{List<String>}, \texttt{Optional<T>}, \ldots) so the compiler can enforce type safety and eliminate many casts. They apply to reference types; primitives use wrappers in generic positions.
\end{solution}

\question[4]
What should be the output for this code snippet?
\begin{lstlisting}[style=java]
import java.util.Collection;
import java.util.HashSet;
import java.util.Set;
import java.util.TreeSet;
public class SortSet {
    public static void main(String...a){
        Collection<Integer> collection = new HashSet<Integer>();
        collection.add(3);
        collection.add(1);
        collection.add(2);
        Set<Integer> treeSet = new TreeSet<Integer>(collection);
        System.out.println(treeSet);
    }
}
\end{lstlisting}
\begin{choices}
  \CorrectChoice \texttt{[1,2,3]}
  \choice \texttt{ClassCastException}
  \choice \texttt{[3,2,1]}
  \choice Compile time error
\end{choices}
\begin{solution}
  \texttt{HashSet} insertion order is not the answer here: the set is passed to the \texttt{TreeSet} constructor. \texttt{TreeSet} sorts elements in \textbf{natural order} (for \texttt{Integer}, ascending), so printing yields \texttt{[1, 2, 3]}.
\end{solution}

\question[4]
Pick the correct definition for \texttt{Queue}.
\begin{choices}
  \CorrectChoice A \texttt{Collection} interface that holds \& processes elements in FIFO (first in, first out) order.
  \choice A \texttt{Collection} interface that holds \& processes elements in FILO (first in, last out) order.
  \choice A \texttt{Collection} interface that holds \& processes elements in LIFO (last in, first out) order.
  \choice A \texttt{Collection} interface that holds \& processes elements in FIRO (first in, random out) order.
\end{choices}
\begin{solution}
  A queue is \textbf{FIFO}: enqueue at one end, dequeue from the other (conceptually). LIFO describes a \textbf{stack}. FILO matches stack behavior; FIRO is not a standard term for queues.
\end{solution}

\question[4]
What is the difference between \texttt{ArrayDeque} and \texttt{PriorityQueue}?
\begin{choices}
  \CorrectChoice \texttt{ArrayDeque} is an implementation of a pure double-ended queue (elements can be added or removed from either end of the queue). \texttt{PriorityQueue} is an implementation of a queue that does not have specified ends (addition or removal of elements doesn't happen on any specific end of the queue).
  \choice \texttt{ArrayDeque} is an implementation of a queue that does not have specified ends (addition or removal of elements doesn't happen on any specific end of the queue). \texttt{PriorityQueue} is an implementation of a pure double-ended queue (elements can be added or removed from either end of the queue).
  \choice \texttt{ArrayDeque} is an implementation of a pure single-ended queue (elements can be added or removed from only one end of the queue). \texttt{PriorityQueue} is an implementation of a queue that does have specified ends (addition or removal of elements can happen on any specific end of the queue).
  \choice \texttt{ArrayDeque} is an implementation of a queue that does have specified ends (addition or removal of elements can happen on any specific end of the queue). \texttt{PriorityQueue} is an implementation of a pure single-ended queue (elements can be added or removed from only one end of the queue).
\end{choices}
\begin{solution}
  \texttt{ArrayDeque} is an efficient \textbf{deque}: \texttt{addFirst}/\texttt{addLast}, 
  \texttt{removeFirst}/\texttt{removeLast}. On the other hand, 
  \texttt{PriorityQueue} is heap-based: the ``next'' element is the 
  \textbf{highest (or lowest) priority} by comparator or natural order, 
  not strictly front or back in a spatial sense.
\end{solution}

\newpage

\question[4]
In linked list implementation of queue, if only front pointer is maintained, which of the following operation take worst case linear time?
\begin{choices}
  \choice Insertion
  \choice Deletion
  \choice To empty a queue
  \CorrectChoice Both Insertion and To empty a queue
\end{choices}
\begin{solution}
  With only a \texttt{front} pointer, enqueue at the tail requires walking the whole list to find the last node---\textbf{O(n)}. Dequeue at the front is \textbf{O(1)}. Emptying the queue by repeated dequeues is \textbf{O(n)} total (linear in size). The option pairing insertion and emptying matches the intended worst-case linear characterization for this design.
\end{solution}

\question[4]
Which of the following methods is highly time efficient? Both methods return fibonacci series.

\noindent Code a:
\begin{lstlisting}[style=java]
int Fib(int n)
{
    if (n <= 1)
        return n;

    return recurssion(n - 1) + recurssion(n - 2);
}
\end{lstlisting}

\noindent Code b:
\begin{lstlisting}[style=java]
int Fib(int n)
{
    int f[] = new int[n + 1];
    f[0] = 0;
    f[1] = 1;
    for (int i = 2; i <= n; i++)
        f[i] = f[i - 1] + f[i - 2];
    return f[n];
}
\end{lstlisting}
\begin{choices}
  \choice Method a
  \CorrectChoice Method b
  \choice Both method a and b
  \choice none of the above
\end{choices}
\begin{solution}
  Code~a (naive recursion without memoization) recomputes the same subproblems exponentially---time roughly exponential in $n$. Code~b fills an array in one pass: \textbf{O(n)} time and \textbf{O(n)} extra space. (The identifier \texttt{recurssion} stands in for a recursive call in the original prompt.)
\end{solution}

\question[4]
Divide and conquer algorithms use recursion.
\begin{choices}
  \CorrectChoice TRUE
  \choice FALSE
\end{choices}
\begin{solution}
  Classic divide and conquer (mergesort, binary search structure, etc.) \textbf{expresses} the pattern by splitting the problem, solving subproblems (often recursively), and combining. Iterative versions exist, but the paradigm is closely tied to recursive formulation.
\end{solution}

\question[4]
Time Complexity. What is time complexity? (choose all that apply)
\begin{checkboxes}
  \choice The actual amount of time it takes an algorithm to execute
  \CorrectChoice A representation of the number of times a block of code need be executed.
  \choice A representation of the amount of memory required to execute an algorithm.
  \CorrectChoice A represntation of how the performance of an algorithm changes as the size of its input scales.
\end{checkboxes}
\begin{solution}
  Time complexity abstracts away wall-clock time and machine details. One view is \textbf{how many times} a dominant operation or block runs as a function of input size $n$. Another (asymptotic) view is \textbf{how running time grows} when $n$ increases---the idea behind Big-O and related notation. The first keyed line matches the step-count view; the second (typo \textit{represntation} preserved) matches the scaling view. Clock time is not the definition; memory growth is \textbf{space} complexity.
\end{solution}

\question[4]
What is the time complexity of a linear search?
\begin{choices}
  \choice $O(1)$
  \CorrectChoice $O(n)$
  \choice $O(\log n)$
  \choice $O(n^2)$
\end{choices}
\begin{solution}
  Scanning an unsorted array or list element by element may examine all $n$ items in the worst case, hence \textbf{linear} time $O(n)$.
\end{solution}

\question[4]
Big-O notation is an example of asymptotic notation that represents the worst case performance for an algorithm.
\begin{choices}
  \CorrectChoice TRUE
  \choice FALSE
\end{choices}
\begin{solution}
  Big-O upper-bounds growth; it is standard to state \textbf{worst-case} $O(f(n))$ unless a problem specifies average or best case. (Other notations: $\Omega$ for lower bounds, $\Theta$ for tight bounds.)
\end{solution}

\question[4]
What is the best case time complexity of bubble sort?
\begin{choices}
  \choice $O(n^2)$
  \choice $O(1)$
  \CorrectChoice $O(n)$
  \choice $O(\log n)$
\end{choices}
\begin{solution}
  With an implementation that detects \textbf{no swaps} on a pass and stops early, an already-sorted array is handled in one pass---\textbf{$O(n)$} comparisons. Without that optimization, best case is still often stated as $O(n^2)$; this question assumes the adaptive variant.
\end{solution}

\question[4]
What is the time complexity of merge sort?
\begin{choices}
  \choice $O(\log n)$
  \CorrectChoice $O(n \log n)$
  \choice $O(n^2)$
  \choice $O(n)$
\end{choices}
\begin{solution}
  Merge sort divides the array in half recursively ($O(\log n)$ levels) and merges linearly at each level, giving \textbf{$O(n \log n)$} in the worst case---a hallmark of efficient comparison sorting.
\end{solution}

\question[4]
In linked list implementation of a queue, front and rear pointers are tracked. Which of these pointers will change during an insertion into a \textbf{NONEMPTY} queue?
\begin{choices}
  \choice Only front pointer
  \choice No pointer will be changed
  \CorrectChoice Only rear pointer
  \choice Both front and rear pointer
\end{choices}
\begin{solution}
  Enqueue links a new node after the current tail and updates \textbf{rear}. \textbf{Front} still points at the first element to be dequeued. For the \textbf{first} insertion into an empty queue, both front and rear typically move; the question restricts to a nonempty queue.
\end{solution}

\question[4]
Entries in a stack are ``ordered''. What is the meaning of this statement?
\begin{choices}
  \CorrectChoice Elements in a stack are added one after another in a specific sequence.
  \choice A collection of stacks is sortable
  \choice Stack entries may be compared with the \texttt{<} operation
  \choice The entries are stored in a linked list
\end{choices}
\begin{solution}
  ``Ordered'' here means there is a \textbf{definite sequence} determined by the order of pushes (LIFO when you pop): later pushes sit above earlier ones. That is not about numeric comparison (\texttt{<}), not about sorting collections of stacks, and not tied to linked lists versus arrays---an implementation detail.
\end{solution}

\newpage

\uplevel{\par\textbf{Lists, maps, sets, and queues (practice bank)}\par}

\question[4]
\textit{(Select all that apply.)} What are some differences between a \texttt{LinkedList} and an \texttt{ArrayList}?
\begin{checkboxes}
  \choice A \texttt{LinkedList} and an \texttt{ArrayList} both only implement the \texttt{List} interface
  \choice An \texttt{ArrayList} is the same as a \texttt{LinkedList}
  \CorrectChoice A \texttt{LinkedList} has nodes with pointers to the next node
  \CorrectChoice A \texttt{LinkedList} has faster insertion and deletion time than an \texttt{ArrayList}
\end{checkboxes}
\begin{solution}
  Both implement \texttt{List} (and \texttt{ArrayList} also implements \texttt{RandomAccess}); they are not the same structure. Linked lists use node chaining; insert/delete in the middle is often cheaper than shifting a large backing array, though big-O depends on how you locate the index.
\end{solution}

\question[4]
The \texttt{Map} interface defines an \texttt{iterator()} method.
\begin{choices}
  \choice TRUE
  \CorrectChoice FALSE
\end{choices}
\begin{solution}
  \texttt{Map} is not \texttt{Iterable}; you iterate \texttt{keySet()}, \texttt{values()}, or \texttt{entrySet()}.
\end{solution}

\question[4]
Pick the correct definition for \texttt{Queue}:
\begin{choices}
  \CorrectChoice A \texttt{Collection} that holds and processes elements in FIFO (first in, first out) order
  \choice A \texttt{Collection} that holds and processes elements in FILO (first in, last out) order
  \choice A \texttt{Collection} that holds and processes elements in LIFO (last in, first out) order
  \choice A \texttt{Collection} that holds and processes elements in LILO (last in, last out) order
\end{choices}
\begin{solution}
  A queue is \textbf{FIFO}; a stack is \textbf{LIFO}. FILO matches a stack, not a queue.
\end{solution}

\question[4]
If you add three times the \emph{same} object reference to a \texttt{HashSet}, how many entries are stored?
\begin{choices}
  \CorrectChoice 1
  \choice 2
  \choice 3
  \choice None of the above
\end{choices}
\begin{solution}
  A set keeps unique elements per \texttt{equals}/\texttt{hashCode}; the same reference is duplicate and stored once.
\end{solution}

\question[4]
\texttt{HashSet} is sorted.
\begin{choices}
  \choice TRUE
  \CorrectChoice FALSE
\end{choices}
\begin{solution}
  \texttt{HashSet} does not define iteration order (unless using a sorted variant such as \texttt{TreeSet}).
\end{solution}

\endgradingrange{csadsa}
